\subsection{Historical supplementary-data reanalysis}

For a package-based reanalysis of the large-scale Table~1 analysis in
\citet{chekouo2017}, we converted the Biometrics supplementary files into the
input format used by \pkg{IntegMultiReg}.  The converted data reproduce the
authors' availability subgroups exactly: \code{111} (E1, all three platforms,
$n=128$), \code{011} (E2, mRNA and miRNA, $n=70$), \code{101} (E3, mRNA and
methylation, $n=120$), and \code{001} (E5, mRNA only, $n=130$).
Table~\ref{tab:kirc-supplement} preserves a historical raw-time run with 400000 retained
MCMC draws, 50000 burn-in draws, and 10 rounds of 10-fold post-fitting
cross-validation.  The same MCMC length and random seed are used for each
IMR/BMS fit. This table predates the correction of survival preprocessing in
version 0.1.3: the old package used raw time, whereas the original supplementary
program takes its logarithm. The table is retained solely as an audit record;
it does not validate the corrected log-time AFT model. The optional replication
path sets \code{survival\_scale = "identity"} explicitly. The long-chain
log-time analysis has not yet been rerun.

\begin{table}[H]
\centering
\small
\begin{tabular}{l c c c c c}
\hline
Model & E1 & E2 & E3 & E5 & Full \\
\hline
CPH + C     & 0.726 (0.008) & 0.719 (0.017) & 0.758 (0.007) & 0.682 (0.015) & 0.719 (0.003) \\
IMR + C + M & 0.869 (0.012) & 0.768 (0.017) & 0.877 (0.007) & 0.782 (0.015) & 0.847 (0.007) \\
BMS + C + M & 0.868 (0.008) & 0.758 (0.025) & 0.865 (0.010) & 0.796 (0.021) & 0.842 (0.007) \\
IMR + M     & 0.853 (0.013) & 0.740 (0.030) & 0.884 (0.012) & 0.699 (0.021) & 0.823 (0.009) \\
BMS + M     & 0.865 (0.015) & 0.714 (0.028) & 0.824 (0.015) & 0.740 (0.017) & 0.817 (0.006) \\
\hline
\end{tabular}
\caption{\label{tab:kirc-supplement}C-index results using the Biometrics
supplementary KIRC data files in the legacy raw-time package implementation.
These are historical results, not results for the corrected log-time model. Parenthetic values are sample standard
deviations across ten post-fitting cross-validation rounds. Three decimal
places are shown. Here C denotes clinical covariates, and M denotes molecular
markers.}
\end{table}
\FloatBarrier

This package-based reanalysis is not a literal duplication of the original
study-specific \proglang{C} program.  In particular, \pkg{IntegMultiReg}
constructs the cross-validation partitions separately within each availability
subgroup, and the order in which model parameters are sampled differs from the
original \proglang{C} implementation. These deliberate design choices are preserved. The additional
raw-time preprocessing discrepancy is a separate implementation issue and
precludes interpreting this historical table as a validation of the intended
log-time model.

The following two subsections record the likelihoods used for the outcome types that
extend the original right-censored AFT formulation of \citet{chekouo2017}.  The
key point is that both extensions can be written in terms of the same latent
Gaussian linear predictor used in Equation~\ref{eq:linear}.  Hence, after
conditioning on the observed or augmented latent response, the pMOM prior, the
MRF prior on the selection indicators and the MCMC updates for
$\bm{\gamma}$ and $\bm{\Theta}$ remain unchanged.

